% \documentclass[10pt,preprint]{sigplanconf}
\documentclass[10pt,twocolumn]{article}

% \sloppy  % avoid overflowing lines ....

\RequirePackage{xspace}
\usepackage[usenames,dvipsnames]{xcolor}
\usepackage{tabularx}
\usepackage{placeins}
\usepackage{multirow} 
\usepackage{makeidx}  % allows for indexgeneration
\usepackage[pdftex]{graphicx} % PNGs
\usepackage{stmaryrd}
\usepackage{bcprules}
% \bcprulessavespacetrue
\bcprulestwocoltrue
\smallrulenamestrue
\setlength{\labelminsep}{1em}

\usepackage{amsmath, amssymb} % algebra
\usepackage[english]{babel} % deutsch
% \usepackage[ngerman]{babel} % deutsch
\usepackage[utf8x]{inputenc}
\usepackage[T1]{fontenc} 
%DIF 27a27
\usepackage{relsize} %DIF > 
%DIF -------
\usepackage[final]{listings} % for sourcecode
\usepackage{graphviz} % graphs
\usepackage{array} % tables
\usepackage{afterpage} % figures
\usepackage{float} % figures
\usepackage{paralist} % figures
\usepackage[shortcuts]{extdash} % figures
\usepackage{hyperref}
\usepackage[textwidth=5em,obeyDraft]{todonotes}

\setlength{\marginparwidth}{1.9cm}
\usepackage{xfrac}

\newcommand{\maybe}{{\bf\footnotesize{?}}\kern-1pt}

\newcommand{\stay}[2]{stay({\tt #1})~$=$~{\tt #2}}
\newcommand{\istay}[2]{i-stay({\tt #1})~$=$~{\tt #2}}
\newcommand{\evalH}[4]{{\tt #1;#2}\,$\vdash$\,{\tt #3}\,$\Downarrow$\,{\tt #4}}
\newcommand{\solveH}[3]{{\tt #1;#2}\,$\models$\,{\tt #3}}
\newcommand{\failH}[1]{$\not\models_{E,H}$\,{\tt #1}}
\newcommand{\stepH}[9]{{\tt <#1|#2|#3|#4|#5>}\,$\longrightarrow$\,{\tt <#6|#7|#8|#9>}}
\newcommand{\tcH}[4]{{\tt #1;#2}\,$\vdash${\tt #3} {\upshape :} {\tt #4}}
\newcommand{\tcHC}[3]{{\tt #1;#2}\,$\vdash$\,{\tt #3}}
\newcommand{\toSolverH}[2]{{\tt E,H}\,$\vdash$\,{\tt #1}\,$\leadsto$\,$#2$}
\newcommand{\fromSolverH}[3]{$#1$\,$\leadsto$\,{\tt #2,#3}}

\newcommand{\inferrule}[3]{\infrule[#1]{\mbox{#2}}{\mbox{#3}}}
\newcommand{\inferax}[2]{\infax[#1]{\mbox{#2}}}

\newcommand{\judgment}[1]{\begin{flushleft}\fbox{#1}\end{flushleft}}

\newcommand{\combine}[4]{{\tt #1;#2}\,$\vdash$\,{\tt #3}\,$\leadsto$\,{\tt #4}}
\newcommand{\inline}[9]{{\tt <#1,#2,#3,#4,#5,#6>}\,$\leadsto$\,{\tt <#7,#8,#9>}}

\newcommand{\E}{\texttt{\upshape E}\xspace}
\newcommand{\Scope}{\texttt{\upshape S}\xspace}
\newcommand{\C}{\texttt{C}}
\newcommand{\I}{\texttt{I}}
\renewcommand{\H}{\texttt{\upshape H}\xspace}
\xspaceaddexceptions{'$'$ |(}
\newcommand{\reinline}[6]{{\tt <#1,#2,#3,#4>} $\leadsto$ {\tt <#5,#6>}}

\def\tdm#1{\todo[size=\footnotesize,color=blue!40]{\textsf{TDM: #1}}}
\def\ab#1{\todo[size=\footnotesize,color=red!40]{\textsf{AB: #1}}}
\def\tf#1{\todo[size=\footnotesize,color=purple!40]{\textsf{TF: #1}}}

\lstdefinelanguage{JavaScript}{
  keywords={typeof, new, true, false, catch, return, null, catch, switch, var, if, in, while, do, else, case, break, this, throw},
  keywordstyle=\bfseries,
  ndkeywords={function, class, export, boolean},
  ndkeywordstyle=\color{darkgray}\bfseries,
  identifierstyle=\color{black},
  sensitive=false,
  comment=[l]{//},
  morecomment=[s]{/*}{*/},
  commentstyle=\color{Bittersweet},%\ttfamily,
  stringstyle=\color{Green},%\ttfamily,
  morestring=[b]',
  morestring=[b]"
}

\lstset{%
    language=JavaScript,%
    basicstyle=\small,%\ttfamily,
    frame=none,%
    numbers=left,
    numbersep=5pt,
    numberstyle=\tiny\color{gray},
    emph={always, once, constructor, ro},%
    emphstyle=\bfseries,%
    % breaklines=true,
    breakatwhitespace=false,%
    literate={<<}{\tiny{$\ll$}~}{1}%
    {?}{\bfseries ?}{1}%
    {priority:}{priority: }{1}%
    {credit:}{credit: }{1}%
    {debit:}{debit: }{1}%
    % {==}{=\slightspace{}=}{1}%
    {visible?}{visible?}{1},%
    showstringspaces=false,%
    columns=fullflexible,%
    frame=b,%
    rulecolor=\color{lightgray},
    % xleftmargin=\parindent,%
}
\lstnewenvironment{code}[1][]%
{\vskip0em\noindent
  \minipage{\linewidth}
  \vspace{0.5\baselineskip}
  \lstset{%
    language=JavaScript,%
    basicstyle=\small%,\ttfamily,
    frame=none,%
    numbers=left,
    numbersep=5pt,
    numberstyle=\tiny\color{gray},
    emph={always, once, constructor, ro},%
    emphstyle=\bfseries,%
    % breaklines=true,
    breakatwhitespace=false,%
    literate={<<}{\tiny{$\ll$}~}{1}%
    {.?}{\bfseries\emph{.?}}{1}%
    % {===}{=\slightspace{}=\slightspace{}=}{1}%
    % {==}{=\slightspace{}=}{1}%
    {visible?}{visible?}{1},%
    showstringspaces=false,%
    columns=flexible,%
    #1}}
{\endminipage\vskip0em}

\restylefloat{figure}

\usepackage[acronym]{glossaries}
\newcommand{\ac}[1]{\gls*{#1}}
\newcommand{\acp}[1]{\glspl*{#1}}
\newcommand{\Acp}[1]{\Gls*{#1}}

% Acronyms
\newacronym{csp}{\textsc{csp}}{Constraint Satisfaction Problem}
\newacronym{cip}{\textsc{cip}}{Constraint Imperative Programming}
\newacronym{api}{\textsc{api}}{application programming interface}
\newacronym{dsl}{\textsc{dsl}}{domain specific language}
\newacronym{vm}{\textsc{vm}}{Virtual Machine}
\newacronym{dom}{\textsc{dom}}{Document Object Model}
\newacronym{ocp}{\textsc{ocp}}{Object-Constraint Programming}
\newacronym{ssa}{\textsc{SSA}}{static-single-assignment}
\newacronym{smt}{\textsc{smt}}{Satisfiability Modulo Theory}
\newacronym{rml}{\textsc{RML}}{the Relational Meta-Language}

\hyphenation{LivelyKernel}
\hyphenation{work-around}

\newcommand\slightspacewide{\kern1em}
\newcommand\slightspace{\kern0.5em}
\newcommand\hairspace{\ifmmode\mskip1mu\else\kern0.08em\fi}
\newcommand{\eg}{e.\hairspace{}g.,\xspace}
\newcommand{\Eg}{E.\hairspace{}g.,\xspace}
\newcommand{\Ie}{I.\hairspace{}e.,\xspace}
\newcommand{\ie}{i.\hairspace{}e.,\xspace}
\newcommand{\cf}[1]{cf.\ {#1}\xspace}
% \newcommand{\todo}[1]{\marginpar{\textsc{todo} #1}}

\newcommand{\langname}{\mbox{Babelsberg}\xspace}
\newcommand{\load}{\textsc{load}\xspace}
\newcommand{\store}{\textsc{store}\xspace}
\newcommand{\javascript}{JavaScript\xspace}
\newcommand{\rml}{\ac{rml}\xspace}
% \newcommand{\rml}{\textsc{RML}\xspace}


\newtheorem{Theorem}{Theorem}
\newtheorem{Corollary}{Corollary}
\newtheorem{Lemma}{Lemma}
\newtheorem{Claim}{Claim}
\newtheorem{Proposition}{Proposition}
\newtheorem{Definition}{Definition}

\newenvironment{theorem}{\noindent\textbf{Theorem.~~}}
{\hfill\\}


%% Commands for one-page artifact description appendices in lncs volumes
% Written by Camil Demetrescu and Erik Ernst
% April 8, 2014

% ARTIFACT: This entire file should be used as-is; it defines standard
% headings to be included in the artifact description, and it will be 'input'
% into the file artifact.tex such that you can use the environments defined
% below

\newcommand{\artifactauthors}[1]{{\paragraph{\bf Authors of the artifact.}#1}}
\newenvironment{summary}{\paragraph{\bf Summary.}}{}
\newenvironment{content}{\paragraph{\bf Content.}}{}
\newenvironment{getting}{\paragraph{\bf Getting the artifact.} }{}
\newenvironment{platforms}{\paragraph{\bf Tested platforms.}}{}
\newcommand{\license}[1]{{\paragraph{\bf License.}#1}}
\newcommand{\mdsum}[1]{{\paragraph{\bf MD5 sum of the artifact.}#1}}
\newcommand{\artifactsize}[1]{{\paragraph{\bf Size of the artifact.}#1}}
\newcommand{\AEBadge}{%
  \raisebox{-9mm}[-15mm][0pt]{%
    \hspace{-39mm}%
    \includegraphics[width=15mm]{aec-badge-ecoop}%
    \hspace{22.5mm}}}

% \usepackage{draftwatermark}
% \SetWatermarkText{\textsf{DRAFT}}
% \SetWatermarkLightness{0.8}
% \SetWatermarkScale{3} 
%DIF PREAMBLE EXTENSION ADDED BY LATEXDIFF
%DIF UNDERLINE PREAMBLE %DIF PREAMBLE
\RequirePackage[normalem]{ulem} %DIF PREAMBLE
\RequirePackage{color}\definecolor{RED}{rgb}{1,0,0}\definecolor{BLUE}{rgb}{0,0,1} %DIF PREAMBLE
\providecommand{\DIFaddtex}[1]{{\protect\color{blue}\uwave{#1}}} %DIF PREAMBLE
\providecommand{\DIFdeltex}[1]{{\protect\color{red}\sout{#1}}}                      %DIF PREAMBLE
%DIF SAFE PREAMBLE %DIF PREAMBLE
\providecommand{\DIFaddbegin}{} %DIF PREAMBLE
\providecommand{\DIFaddend}{} %DIF PREAMBLE
\providecommand{\DIFdelbegin}{} %DIF PREAMBLE
\providecommand{\DIFdelend}{} %DIF PREAMBLE
%DIF FLOATSAFE PREAMBLE %DIF PREAMBLE
\providecommand{\DIFaddFL}[1]{\DIFadd{#1}} %DIF PREAMBLE
\providecommand{\DIFdelFL}[1]{\DIFdel{#1}} %DIF PREAMBLE
\providecommand{\DIFaddbeginFL}{} %DIF PREAMBLE
\providecommand{\DIFaddendFL}{} %DIF PREAMBLE
\providecommand{\DIFdelbeginFL}{} %DIF PREAMBLE
\providecommand{\DIFdelendFL}{} %DIF PREAMBLE
%DIF END PREAMBLE EXTENSION ADDED BY LATEXDIFF
%DIF PREAMBLE EXTENSION ADDED BY LATEXDIFF
%DIF HYPERREF PREAMBLE %DIF PREAMBLE
\providecommand{\DIFadd}[1]{\texorpdfstring{\DIFaddtex{#1}}{#1}} %DIF PREAMBLE
\providecommand{\DIFdel}[1]{\texorpdfstring{\DIFdeltex{#1}}{}} %DIF PREAMBLE
%DIF END PREAMBLE EXTENSION ADDED BY LATEXDIFF



\begin{document}

\special{papersize=8.5in,11in}
\setlength{\pdfpageheight}{\paperheight}
\setlength{\pdfpagewidth}{\paperwidth}
\setlength{\hoffset}{-1in}
\setlength{\voffset}{-1in}
\setlength{\topmargin}{1in}
\setlength{\headheight}{0pt}
\setlength{\headsep}{0pt}
\setlength{\evensidemargin}{.75in}
\setlength{\oddsidemargin}{.75in}
\newdimen{\standardtextwidth}
\setlength{\standardtextwidth}{42pc}
\setlength{\textwidth}{\standardtextwidth}
\setlength{\topskip}{8pt}
\setlength{\columnsep}{2pc}
\setlength{\textheight}{54.5pc}
\setlength{\parskip}{0pt}
\setlength{\parindent}{12pt}
\setlength{\lineskip}{.5pt}
\setlength{\lineskiplimit}{\lineskip}
\frenchspacing
\pretolerance = 400
\tolerance = \pretolerance
\setlength{\emergencystretch}{5pt}
\clubpenalty = 10000
\widowpenalty = 10000
\setlength{\hfuzz}{.5pt}
\setlength{\marginparwidth}{36pt}
\setlength{\marginparsep}{2pt}
\setlength{\marginparpush}{8pt}
\setlength{\skip\footins}{8pt plus 3pt minus 1pt}
\setlength{\footnotesep}{9pt}

\title{Cover Letter for Revisions to OOPSLA'15 Paper 211
\\\large Checks and Balances: Constraint Solving without Surprises in Object-Constraint Programming Languages}
\date{}

% Communications Design Group, SAP Labs, San Francisco, CA, USA}

\maketitle

\paragraph{Connection between constraint variables and ordinary variables}
Regarding this issue and how this affects what the solver can do with
respect to adding/removing fields, we have made an effort to make this
clearer. First, we added the key property that all variables can be
used in constraints and potentially solved for when first introducing
the example in Listing 1:
\begin{quote}
The constraint on line 2 ensures that \verb!acc.balance! remains
\DIFdelbegin \DIFdel{at least }%DIFDELCMD < \verb!minBal!%%%
\DIFdel{.  }\DIFdelend \DIFaddbegin \DIFadd{above the minimum balance
  required by the plan.  Note that the plan's minimum balance is the
  result of a message send---it does not matter if it is calculated or
  a field access. }\DIFaddend If the constraint is already satisfied
at that point, then nothing need be done; otherwise the solver will
modify the values of variables and their \DIFdelbegin
\DIFdel{properties }\DIFdelend \DIFaddbegin \DIFadd{fields }\DIFaddend
in order to satisfy the constraint. \DIFaddbegin \DIFadd{Note also
  that these variables and fields are not special ``constraint
  variables''---any variable can appear and potentially be solved for
  in a constraint. }\DIFaddend  
\end{quote}
We think some of the confusion may have also been due to the fact
that, in the semantics, we simply replace the entire heap and
environment with those generated by the solver - this is just for
simplicity and really an implementation detail, so we removed
this. Instead, we have attempted to clarify that the issue really
arises when we allow any kind of constraints without type-checking
them. A relevant text was added in Section 2.1:
\begin{quote}
Specifically, suppose the object \verb!acc! has no \verb!balance!
field.  In earlier versions of Babelsberg\DIFaddbegin \DIFadd{, }\DIFaddend the solver could silently
invent such a field \DIFdelbegin \DIFdel{and then constrain it. }\DIFdelend \DIFaddbegin \DIFadd{to satisfy the constraint on line 2, and that field
would be added to the object when the solution is used to update the
heap---after all, we did ask in our constraint that such a field exist
and be equal to the minimum balance. (Similarly, if we created a
constraint that the account be equal to an object with only a balance
field and if the account previously had additional fields, the superfluous
fields could be removed.) 
}\DIFaddend While one might argue that such behavior is desirable in \DIFdelbegin \DIFdel{this case}\DIFdelend \DIFaddbegin \DIFadd{these
cases}\DIFaddend , it is a slippery slope.  For example, consider a constraint of
the form \verb!p.x > 0 || p.y > 0!  on a point \verb!p!.  There is
inherent nondeterminism here --- the solver can choose which \DIFdelbegin \DIFdel{property }\DIFdelend \DIFaddbegin \DIFadd{field }\DIFaddend to
update if the constraint does not hold.  But if \verb!p! lacks an
\verb!x!  \DIFdelbegin \DIFdel{property }\DIFdelend \DIFaddbegin \DIFadd{field }\DIFaddend should it be invented?  If \verb!p! lacks both \DIFdelbegin \DIFdel{properties }\DIFdelend \DIFaddbegin \DIFadd{fields
}\DIFaddend which one should be invented?
\end{quote}

\paragraph{Lifetime of always constraints}
\label{sec:lifet-always-constr}
Regarding the asymmetry of always constraints only applying from the
point at which they are executed (but not retroactively), and
semantically persisting forever (implying interaction with automatic
garbage collection), we have added a clarifications around Listing 1:
\begin{quote}
An \verb!always! constraint is required to
hold for the rest of the execution; Babelsberg/JS therefore \DIFaddbegin \DIFadd{re-creates
and }\DIFaddend re-solves the constraint as part of the execution of each
subsequent statement if it may have become invalidated.  In contrast,
a \verb!once! constraint is solved once and then immediately
\DIFdelbegin \DIFdel{redacted. }\DIFdelend \DIFaddbegin \DIFadd{retracted. Semantically, }\verb!always!  \DIFadd{constraints persist for the
rest of the execution, but the practical implementations garbage
collect constraints if all variables that they transitively refer to are no longer
accessible.
}\DIFaddend 
\end{quote}
\begin{center}
[\dots]
\end{center}
\begin{quote}
  Because this is an \verb!always! constraint, it will be re-checked,
  and if necessary re-solved, \DIFdelbegin \DIFdel{after }\DIFdelend
  \DIFaddbegin \DIFadd{as part of executing }\DIFaddend each
  subsequent statement \DIFdelbegin \DIFdel{.  It is }\DIFdelend
  \DIFaddbegin \DIFadd{as well.$^1$}\vskip0.3\baselineskip
  {\footnotesize \DIFadd{$^1$The constraint is active from the point
      at which it is created and onwards, but does not retroactively
      apply prior to the execution.  In a graphical application, for
      example, the user may have already seen previous results.}}
  \DIFaddend
\end{quote}

\paragraph{Distinct identity constraints}

There was a question of whether it is possible to express that two
variables must not refer to the same object. It is not, because this
leads to identity non-determinism:

\begin{quote}
  Finally, the solver \DIFdelbegin \DIFdel{can }\DIFdelend
  \DIFaddbegin \DIFadd{could }\DIFaddend also satisfy the identity
  constraint by assigning both variables to yet a third
  object\DIFaddbegin \DIFadd{$^1$}
  \vskip0.3\baselineskip
  {\footnotesize \DIFadd{$^1$One could also imagine a
    }\emph{\DIFadd{distinct}} \DIFadd{identity constraint to express
      that two variables should not refer to the same
      object---however, this will almost always lead to
      non-determinism as the solver is free to pick any object to
      re-satisfy the constraint if one of the variables is assigned to
      be identical to the other.}} \DIFaddend
\end{quote}

\paragraph{Identity preservation and changing references}

``Identity preservation: The utility of this feature tends to
blur. Basically, it means that only values (integers, Booleans, etc.)
can change''

``assignment to reference-type fields are outlawed as well''

`` 'first .. solves identity constaints': But that could never happen,
because the solver is not allowed to change the value of a
reference-typed slot''

\paragraph{Structural determinism}

``What is the structure of an object? For instance, does it include
the length of an array?''

\paragraph{Two-phase solving}

``For instance, Smalltalk, Dart, and others do not have primitive
values, or at least only as an optimized representation of something
which is conceptually and observationally an object''

\paragraph{Typechecking}

``I wasn't sure if the structural typing is really responsible here; I
think it's the dynamic nature of the checks. A dynamic type system can
work for both languages that have static and dynamic type checking.''

``Can the runtime environment access the source representation? Does
it need to?''

``The (T-*) rules are really not well-formed, using the runtime
environment (E;H). How would we typecheck programs at "compile-time" if the
type rules require a stack and a heap to work? Can we define the stack/heap to
be "empty" (whatever that means for a stack when viewed from a method body)?
Maybe the calculus is intended to be imperative-only, and that's the reason
why it does not include classes or methods, and indeed any kind of procedural
abstraction. Please explain.''

\paragraph{Backtracking}

`` Also, it seems that if constraints depend on certain conditions, it
should be possible to retract them. Specifically, is it possible to
intertwine procedural search (using backtracking) with constraint
solving? This would require that constraints are (automatically)
removed upon backtracking …''

\paragraph{Mozart/Oz}

\paragraph{Background information on Babelsberg and its motivations}

``the paper would benefit from additional (or rather, more focused)
background on BabelsbergJS''

\paragraph{``Changing the identity of objects''}

``I find it extremely confusing when the authors speak of “changing
the identity of objects” identity is immutable for all objects. I
think (but cannot be sure) what the authors mean is “change the value
of reference-typed variables”, but this is something entirely
different.''

\paragraph{Formalism}

``The syntax for statements seems to be missing''

`` The rules are also somewhat repetitive, because certain easy
generalizations have not been used (near-identical rules, for
instance).''

``Note that the rules (E-Op), (E-Compare), (E-Combine) can be reduced to
a single rule using a symbol to stand for any of \{\\plusring, \\lessring,
\\orring \}. Similarly for (E-Const), (E-Ref).''

``On p.7 and following: the definitions of stay and i-stay are central
to this paper, but are only given in English. For completeness, they
could be spelled out in mathematical notation.''

`` On p.9, the sentence explaining the meaning of $L \neq x \implies$
... should appear before S-Asgn, which also uses it.''

`` On p.9, in S-AsgnNew, the third-to-last premise should probably be
E;H |- ...  instead of E';H' |- ''

\paragraph{Misc Presentation}

`` I don't like using "overviews" as a verb form ''

``On p.12, why are there asterisks before the last paragraph?''

\paragraph{One-way functions}

``Actually, you're saying that a pure function (no side-effects, body
simply returns an expression) can be run bi-directionally''

\end{document}


%%% Local Variables:
%%% mode: latex
%%% TeX-master: t
%%% End:
